\section{Conclusion and Future Directions}

We presented \textbf{MemeSonic}, a unified system for affective meme understanding and expressive
audio generation. Our central contribution is the operationalization of \emph{cross-modal
incongruity}—the intentional semantic conflict between meme image and text—as a continuous,
differentiable signal that drives both fusion routing and TTS prosody conditioning.

\textbf{Main findings.}
(1) Explicit incongruity modeling improves classification on non-literal affect tasks by up to
+12.1\% accuracy and +14.5\% Macro-F1 over standard fusion, with the largest gains on
metaphor and sarcasm—exactly the tasks where cross-modal conflict is semantically load-bearing.
(2) Incongruity-conditioned audio generation raises pragmatic appropriateness from AMOS 1.68
(Vanilla TTS) to 3.95, with 84.5\% human preference, demonstrating that the pragmatic signal
latent in image-text incongruity can be successfully translated into expressive prosody.
(3) Audio modality carries non-redundant affective information (KL $V\leftrightarrow A = 1.14$
vs.\ KL $V\leftrightarrow T = 0.69$), but naive fusion fails entirely due to modality-space
incompatibility rather than absence of signal—a distinction with practical implications for the
design of tri-modal systems.

\textbf{Limitations.}

\emph{Scale.}
Our models rely on frozen CLIP ViT-B/32 embeddings and may underperform on tasks requiring
fine-grained cultural knowledge. The tri-modal alignment study covers only 100 memes, so the
+0.55 accuracy gain from the linear projector should be read as a diagnostic rather than a
deployable result. Human audio evaluation is likewise small (12 raters, 84 ratings).

\emph{Data quality and annotation subjectivity.}
Sarcasm and humor benchmarks (e.g., Memotion, MMSD) are not meme-domain specific and reflect
social-media text annotation conventions. Within meme datasets, most prior work targets hate
speech—a comparatively low-subjectivity task—while humor intensity and sarcasm type exhibit
substantial inter-annotator disagreement, consistent with our low F1 on humour (36.2\%) and
sarcasm (33.3\%). No existing benchmark jointly annotates sentiment, humor, sarcasm, and audio
for the same memes, precluding end-to-end tri-modal evaluation.

\textbf{Future directions.}
\emph{(1) Meme-specific dataset construction.}
The most impactful near-term contribution would be a large-scale, meme-native dataset with
multi-dimensional affective annotations—sentiment, humor type, sarcasm mechanism, incongruity
score, and paired voiced audio—collected with rigorous inter-annotator agreement protocols.
Such a resource would enable both training and evaluation of the full MemeSonic pipeline without
relying on heterogeneous benchmarks.

\emph{(2) Generalization to broader incongruity domains.}
The incongruity-aware framework developed here is not inherently meme-specific. Sarcasm and humor
in news captions, social media posts, political cartoons, and video shorts share the same
image-text conflict structure. Extending MemeSonic to these domains—particularly short-form video
where temporal incongruity (e.g., a cheerful musical score against distressing imagery) is
prevalent—would demonstrate the generality of the incongruity signal as an architectural primitive.

\emph{(3) Interactive user studies with human-in-the-loop audio creation.}
Rather than passive preference ratings, future work should engage users as active participants in
meme audio creation. A demo-based study could present the meme alongside a palette of prosodically
distinct audio options—spanning emotional tones, speech rates, and incongruity levels—and ask users
to select, remix, or record their preferred voicing. Aggregating these choices would yield
naturalistic, intent-labeled prosody data far richer than post-hoc MOS ratings. Such a study could
also surface cultural and individual variation in how incongruity is perceived and performed,
motivating personalized prosody adaptation grounded in demonstrated user preference rather than
inferred affective labels.

\emph{(4) Joint tri-modal pre-training.}
Our tri-modal alignment study reveals a fundamental obstacle: ImageBind, despite being trained to
bind image, text, and audio in a shared space, produces near-zero cosine similarity between its
audio embeddings and TTS-synthesized speech (Top-1 accuracy $\approx 0\%$), because its audio
encoder was trained on natural soundscapes rather than expressive speech. Bridging this gap with a
post-hoc linear projector yields only a modest +0.55 accuracy improvement, suggesting that
surface-level alignment is insufficient. A more principled fix is to train a shared encoder on
meme image, text, and voiceover audio simultaneously—using incongruity-weighted contrastive
learning as a pre-training objective—so that the resulting embeddings natively support tri-modal
fusion without relying on incompatible pretrained representations.

\emph{(5) LLM + incongruity as structured input.}
Providing the incongruity score $\delta$ and per-modality emotion distributions as structured
features to an LLM (via function calling or tool use) could combine the commonsense reasoning
advantage of large models with the explicit conflict signal of our architecture, potentially
closing the remaining gap on tasks like intention and sarcasm where LLMs currently outperform
trained models.

MemeSonic demonstrates that meme understanding requires moving beyond the assumption of cross-modal
agreement, and that the \emph{conflict itself} is the signal. We hope this work encourages the
community to treat incongruity not as noise to be averaged away, but as a first-class feature to
be modeled, routed on, and generated from.
